% Options for packages loaded elsewhere
\PassOptionsToPackage{unicode}{hyperref}
\PassOptionsToPackage{hyphens}{url}
\PassOptionsToPackage{dvipsnames,svgnames,x11names}{xcolor}
\documentclass[
  11pt,
]{article}
\pdftrailerid{}
\usepackage{xcolor}
\usepackage{amsmath,amssymb}
\setcounter{secnumdepth}{-\maxdimen} % remove section numbering
\usepackage{iftex}
\ifPDFTeX
  \usepackage[T1]{fontenc}
  \usepackage[utf8]{inputenc}
  \usepackage{textcomp} % provide euro and other symbols
\else % if luatex or xetex
  \usepackage{unicode-math} % this also loads fontspec
  \defaultfontfeatures{Scale=MatchLowercase}
  \defaultfontfeatures[\rmfamily]{Ligatures=TeX,Scale=1}
\fi
\usepackage{lmodern}
\ifPDFTeX\else
  % xetex/luatex font selection
\fi
% Use upquote if available, for straight quotes in verbatim environments
\IfFileExists{upquote.sty}{\usepackage{upquote}}{}
\IfFileExists{microtype.sty}{% use microtype if available
  \usepackage[]{microtype}
  \UseMicrotypeSet[protrusion]{basicmath} % disable protrusion for tt fonts
}{}
\makeatletter
\@ifundefined{KOMAClassName}{% if non-KOMA class
  \IfFileExists{parskip.sty}{%
    \usepackage{parskip}
  }{% else
    \setlength{\parindent}{0pt}
    \setlength{\parskip}{6pt plus 2pt minus 1pt}}
}{% if KOMA class
  \KOMAoptions{parskip=half}}
\makeatother
\usepackage{color}
\usepackage{fancyvrb}
\newcommand{\VerbBar}{|}
\newcommand{\VERB}{\Verb[commandchars=\\\{\}]}
\DefineVerbatimEnvironment{Highlighting}{Verbatim}{commandchars=\\\{\}}
% Add ',fontsize=\small' for more characters per line
\usepackage{framed}
\definecolor{shadecolor}{RGB}{248,248,248}
\newenvironment{Shaded}{\begin{snugshade}}{\end{snugshade}}
\newcommand{\AlertTok}[1]{\textcolor[rgb]{0.94,0.16,0.16}{#1}}
\newcommand{\AnnotationTok}[1]{\textcolor[rgb]{0.56,0.35,0.01}{\textbf{\textit{#1}}}}
\newcommand{\AttributeTok}[1]{\textcolor[rgb]{0.13,0.29,0.53}{#1}}
\newcommand{\BaseNTok}[1]{\textcolor[rgb]{0.00,0.00,0.81}{#1}}
\newcommand{\BuiltInTok}[1]{#1}
\newcommand{\CharTok}[1]{\textcolor[rgb]{0.31,0.60,0.02}{#1}}
\newcommand{\CommentTok}[1]{\textcolor[rgb]{0.56,0.35,0.01}{\textit{#1}}}
\newcommand{\CommentVarTok}[1]{\textcolor[rgb]{0.56,0.35,0.01}{\textbf{\textit{#1}}}}
\newcommand{\ConstantTok}[1]{\textcolor[rgb]{0.56,0.35,0.01}{#1}}
\newcommand{\ControlFlowTok}[1]{\textcolor[rgb]{0.13,0.29,0.53}{\textbf{#1}}}
\newcommand{\DataTypeTok}[1]{\textcolor[rgb]{0.13,0.29,0.53}{#1}}
\newcommand{\DecValTok}[1]{\textcolor[rgb]{0.00,0.00,0.81}{#1}}
\newcommand{\DocumentationTok}[1]{\textcolor[rgb]{0.56,0.35,0.01}{\textbf{\textit{#1}}}}
\newcommand{\ErrorTok}[1]{\textcolor[rgb]{0.64,0.00,0.00}{\textbf{#1}}}
\newcommand{\ExtensionTok}[1]{#1}
\newcommand{\FloatTok}[1]{\textcolor[rgb]{0.00,0.00,0.81}{#1}}
\newcommand{\FunctionTok}[1]{\textcolor[rgb]{0.13,0.29,0.53}{\textbf{#1}}}
\newcommand{\ImportTok}[1]{#1}
\newcommand{\InformationTok}[1]{\textcolor[rgb]{0.56,0.35,0.01}{\textbf{\textit{#1}}}}
\newcommand{\KeywordTok}[1]{\textcolor[rgb]{0.13,0.29,0.53}{\textbf{#1}}}
\newcommand{\NormalTok}[1]{#1}
\newcommand{\OperatorTok}[1]{\textcolor[rgb]{0.81,0.36,0.00}{\textbf{#1}}}
\newcommand{\OtherTok}[1]{\textcolor[rgb]{0.56,0.35,0.01}{#1}}
\newcommand{\PreprocessorTok}[1]{\textcolor[rgb]{0.56,0.35,0.01}{\textit{#1}}}
\newcommand{\RegionMarkerTok}[1]{#1}
\newcommand{\SpecialCharTok}[1]{\textcolor[rgb]{0.81,0.36,0.00}{\textbf{#1}}}
\newcommand{\SpecialStringTok}[1]{\textcolor[rgb]{0.31,0.60,0.02}{#1}}
\newcommand{\StringTok}[1]{\textcolor[rgb]{0.31,0.60,0.02}{#1}}
\newcommand{\VariableTok}[1]{\textcolor[rgb]{0.00,0.00,0.00}{#1}}
\newcommand{\VerbatimStringTok}[1]{\textcolor[rgb]{0.31,0.60,0.02}{#1}}
\newcommand{\WarningTok}[1]{\textcolor[rgb]{0.56,0.35,0.01}{\textbf{\textit{#1}}}}
\setlength{\emergencystretch}{3em} % prevent overfull lines
\providecommand{\tightlist}{%
  \setlength{\itemsep}{0pt}\setlength{\parskip}{0pt}}
% obstruction_preamble.tex — shared house style for ALL papers.
% Engine: pdfLaTeX (portable, no font cache). Injected by pandoc (see makefile).
% pandoc loads fontenc(T1)+inputenc(utf8) before this; hyperref after (so we set
% link colors via pandoc -V flags, not here).

% ---- Fonts & math ----
\IfFileExists{lmodern.sty}{\usepackage{lmodern}}{}   % Latin Modern if present
\usepackage{amsmath,amssymb,amsthm}                  % \mathbb \mathfrak \square ...
\usepackage{microtype}

% ---- Over-wide \boxed kernel statements shrink to fit the text width ----
% A long \boxed{\text{...}} line can exceed \linewidth because \text does not
% line-break in math mode. "max width=\linewidth" scales down ONLY boxes that
% are wider than the text block; short math boxes are left at natural size.
% One place, all papers.
\usepackage{adjustbox}
\let\oboldboxed\boxed
\renewcommand{\boxed}[1]{\adjustbox{max width=\linewidth}{\oboldboxed{#1}}}
\setlength{\emergencystretch}{3em}

% ---- Line breaking for verbatim and code blocks ----
\usepackage{fancyvrb}

% ---- Safe line breaking for typewriter font & long URLs/DOIs ----
\usepackage[htt]{hyphenat}
\usepackage{xurl}
\usepackage[hyphens]{url}


% ---- Flawless fix: literal Unicode (ℓ¹ Φ₁ κ¹ ⊗ →) in TEXT mode ----
% In $...$ these already work; this maps stray text-mode glyphs to math so
% nothing renders blank and you never rewrite source. If a build ever says
% "Unicode character X not set up", add one line here for that glyph.
\usepackage{newunicodechar}
\newcommand{\um}[1]{\ensuremath{#1}}
\newunicodechar{ℓ}{\um{\ell}}
\newunicodechar{Φ}{\um{\Phi}}\newunicodechar{φ}{\um{\varphi}}
\newunicodechar{Ψ}{\um{\Psi}}\newunicodechar{ψ}{\um{\psi}}
\newunicodechar{⇔}{\um{\Leftrightarrow}}
\newunicodechar{Γ}{\um{\Gamma}}\newunicodechar{γ}{\um{\gamma}}
\newunicodechar{Ω}{\um{\Omega}}\newunicodechar{ω}{\um{\omega}}
\newunicodechar{κ}{\um{\kappa}}\newunicodechar{χ}{\um{\chi}}
\newunicodechar{Λ}{\um{\Lambda}}\newunicodechar{λ}{\um{\lambda}}
\newunicodechar{Θ}{\um{\Theta}}\newunicodechar{θ}{\um{\theta}}
\newunicodechar{Δ}{\um{\Delta}}\newunicodechar{δ}{\um{\delta}}
\newunicodechar{Σ}{\um{\Sigma}}\newunicodechar{σ}{\um{\sigma}}
\newunicodechar{Π}{\um{\Pi}}\newunicodechar{π}{\um{\pi}}
\newunicodechar{α}{\um{\alpha}}\newunicodechar{β}{\um{\beta}}
\newunicodechar{ε}{\um{\varepsilon}}\newunicodechar{μ}{\um{\mu}}
\newunicodechar{ρ}{\um{\rho}}\newunicodechar{τ}{\um{\tau}}
\newunicodechar{ν}{\um{\nu}}\newunicodechar{ξ}{\um{\xi}}
\newunicodechar{ζ}{\um{\zeta}}\newunicodechar{η}{\um{\eta}}
\newunicodechar{ℝ}{\um{\mathbb{R}}}\newunicodechar{ℚ}{\um{\mathbb{Q}}}
\newunicodechar{ℤ}{\um{\mathbb{Z}}}\newunicodechar{ℕ}{\um{\mathbb{N}}}
\newunicodechar{ℂ}{\um{\mathbb{C}}}
\newunicodechar{𝒟}{\um{\mathcal{D}}}
\newunicodechar{→}{\um{\to}}\newunicodechar{↦}{\um{\mapsto}}
\newunicodechar{≤}{\um{\leq}}\newunicodechar{≥}{\um{\geq}}
\newunicodechar{≠}{\um{\neq}}\newunicodechar{≈}{\um{\approx}}
\newunicodechar{≅}{\um{\cong}}\newunicodechar{≡}{\um{\equiv}}
\newunicodechar{⊗}{\um{\otimes}}\newunicodechar{⊕}{\um{\oplus}}
\newunicodechar{∘}{\um{\circ}}\newunicodechar{∂}{\um{\partial}}
\newunicodechar{∇}{\um{\nabla}}\newunicodechar{∑}{\um{\sum}}
\newunicodechar{∏}{\um{\prod}}\newunicodechar{∞}{\um{\infty}}
\newunicodechar{⊆}{\um{\subseteq}}\newunicodechar{⊂}{\um{\subset}}
\newunicodechar{⊇}{\um{\supseteq}}\newunicodechar{∈}{\um{\in}}
\newunicodechar{∉}{\um{\notin}}\newunicodechar{∩}{\um{\cap}}
\newunicodechar{∪}{\um{\cup}}\newunicodechar{∖}{\um{\setminus}}
\newunicodechar{∅}{\um{\varnothing}}\newunicodechar{⊥}{\um{\perp}}
\newunicodechar{∀}{\um{\forall}}\newunicodechar{∃}{\um{\exists}}
\newunicodechar{±}{\um{\pm}}\newunicodechar{×}{\um{\times}}\newunicodechar{−}{\um{-}}
\newunicodechar{·}{\um{\cdot}}\newunicodechar{√}{\um{\surd}}
\newunicodechar{⟨}{\um{\langle}}\newunicodechar{⟩}{\um{\rangle}}
\newunicodechar{≪}{\um{\ll}}\newunicodechar{≫}{\um{\gg}}
\newunicodechar{∝}{\um{\propto}}\newunicodechar{≜}{\um{\triangleq}}
\newunicodechar{⇒}{\um{\Rightarrow}}\newunicodechar{⇐}{\um{\Leftarrow}}
\newunicodechar{∧}{\um{\wedge}}\newunicodechar{∨}{\um{\vee}}
\newunicodechar{⟺}{\um{\iff}}\newunicodechar{□}{\um{\square}}
\newunicodechar{⊔}{\um{\sqcup}}\newunicodechar{…}{\dots}\newunicodechar{‖}{\um{\Vert}}
\newunicodechar{▼}{\um{\blacktriangledown}}\newunicodechar{«}{\guillemotleft}\newunicodechar{»}{\guillemotright}
\newunicodechar{─}{--}\newunicodechar{│}{|}\newunicodechar{┌}{+}\newunicodechar{┐}{+}\newunicodechar{┼}{+}
\newunicodechar{═}{=}\newunicodechar{║}{\um{\parallel}}\newunicodechar{̄}{\={}}
\newunicodechar{⟹}{\um{\implies}}\newunicodechar{⟸}{\um{\impliedby}}
\newunicodechar{⁰}{\textsuperscript{0}}\newunicodechar{¹}{\textsuperscript{1}}
\newunicodechar{²}{\textsuperscript{2}}\newunicodechar{³}{\textsuperscript{3}}
\newunicodechar{⁴}{\textsuperscript{4}}\newunicodechar{⁵}{\textsuperscript{5}}
\newunicodechar{⁶}{\textsuperscript{6}}\newunicodechar{⁷}{\textsuperscript{7}}
\newunicodechar{⁸}{\textsuperscript{8}}\newunicodechar{⁹}{\textsuperscript{9}}
\newunicodechar{⁻}{\textsuperscript{-}}\newunicodechar{⁺}{\textsuperscript{+}}
\newunicodechar{ⁿ}{\textsuperscript{n}}
\newunicodechar{₀}{\textsubscript{0}}\newunicodechar{₁}{\textsubscript{1}}\newunicodechar{ᵀ}{\textsuperscript{T}}
\newunicodechar{₂}{\textsubscript{2}}\newunicodechar{₃}{\textsubscript{3}}
\newunicodechar{ₜ}{\textsubscript{t}}\newunicodechar{ₑ}{\textsubscript{e}}\newunicodechar{ₙ}{\textsubscript{n}}

% ---- Geometry & spacing ----
\usepackage{geometry}\geometry{letterpaper,margin=1in}
\setlength{\parindent}{0pt}\setlength{\parskip}{0.55\baselineskip}\linespread{1.04}

% ---- Palette ----
\usepackage{xcolor}
\definecolor{oblink}{HTML}{1A4D8F}\definecolor{obrule}{HTML}{C2CCD6}
\definecolor{obboxbg}{HTML}{F6F8FA}
\definecolor{obproven}{HTML}{1B7F4B}   % PROVEN / DERIVED / STANDARD
\definecolor{obconj}{HTML}{B45309}     % CONJECTURE / OPEN / FALSE
\definecolor{obdef}{HTML}{6B21A8}      % DEFINITION / INTERPRETIVE
\definecolor{obneutral}{HTML}{475569}

% ---- Headings ----
\usepackage{titlesec}
\titleformat{\section}{\Large\bfseries}{\thesection}{0.6em}{}
\titleformat{\subsection}{\large\bfseries}{\thesubsection}{0.6em}{}
\titleformat{\subsubsection}{\normalsize\bfseries}{\thesubsubsection}{0.6em}{}
\titlespacing*{\section}{0pt}{1.4em}{0.5em}
\titlespacing*{\subsection}{0pt}{1.1em}{0.4em}

% ---- Theorem environments (for when you adopt semantic theorems) ----
\theoremstyle{plain}
\newtheorem{theorem}{Theorem}[section]
\newtheorem{lemma}[theorem]{Lemma}
\newtheorem{proposition}[theorem]{Proposition}
\newtheorem{corollary}[theorem]{Corollary}
\newtheorem{conjecture}[theorem]{Conjecture}
\theoremstyle{definition}\newtheorem{definition}[theorem]{Definition}
\theoremstyle{remark}\newtheorem*{remark}{Remark}

% ---- Blockquotes -> one uniform left-ruled callout ----
% pandoc wraps every `>` block (your theorems, Status notes, caveats) in `quote`.
\usepackage{mdframed}
\renewenvironment{quote}{%
  \begin{mdframed}[linewidth=0pt,leftline=true,linecolor=obrule,
    rightline=false,topline=false,bottomline=false,
    innerleftmargin=1em,innerrightmargin=0.8em,innertopmargin=0.5em,
    innerbottommargin=0.5em,skipabove=0.6em,skipbelow=0.6em,
    backgroundcolor=obboxbg]\small}{\end{mdframed}}

% ---- Claim-type tag macro ----
\newcommand{\claimtag}[2]{{\small\sffamily\bfseries\color{#1}[#2]}}

% ---- Inline code / file references -> standard breakable typewriter font ----
\DeclareRobustCommand{\texttt}[1]{%
  \ifmmode
    \mathtt{#1}%
  \else
    {\ttfamily\small\renewcommand{\_}{\symbol{95}\allowbreak}#1}%
  \fi}



% ---- "What this does not claim" box ----
\usepackage[most]{tcolorbox}
\newtcolorbox{nonclaim}{colback=obboxbg,colframe=obrule,boxrule=0.6pt,arc=2pt,
  left=1em,right=1em,top=0.5em,bottom=0.5em,
  title={What this does not claim},fonttitle=\bfseries,
  coltitle=black,colbacktitle=obboxbg}

% ---- Tables ----
\usepackage{booktabs}\renewcommand{\arraystretch}{1.15}
\usepackage{etoolbox}
% Compact tables, and let over-long \mathsf{CamelCase} status identifiers wrap
% inside cells. In math, \mathsf{ContinuumShadowBecomesFakeAuthority} is a single
% unbreakable atom and overruns narrow pandoc columns. Scoped to tables only, we
% re-render \mathsf as a sequence of per-word \mathsf atoms split before each
% interior capital and joined by \allowbreak, so the identifier wraps at the math
% top level (an \allowbreak between atoms IS a legal break in inline math; a text
% \mbox/\seqsplit is not). Also shrink the table font and tighten column padding.
% Body math is untouched; table cells carry only plain ASCII identifiers, so the
% splitter is emoji-safe. One place, all papers.
\let\obmathsforig\mathsf
\ExplSyntaxOn
\tl_new:N \l__obt_run_tl
\int_new:N \l__obt_i_int
\cs_new_protected:Npn \__obt_emit:
  {
    \tl_if_empty:NF \l__obt_run_tl
      {
        \int_compare:nNnT { \l__obt_i_int } > { 1 } { \allowbreak }
        \exp_args:NV \obmathsforig \l__obt_run_tl
        \int_incr:N \l__obt_i_int
        \tl_clear:N \l__obt_run_tl
      }
  }
\cs_new_protected:Npn \obtablesf #1
  {
    \tl_clear:N \l__obt_run_tl
    \int_set:Nn \l__obt_i_int { 1 }
    \tl_map_inline:nn {#1}
      {
        \str_if_eq:eeT { \str_uppercase:n {##1} } {##1}
          { \str_if_eq:eeF { \str_lowercase:n {##1} } {##1} { \__obt_emit: } }
        \tl_put_right:Nn \l__obt_run_tl {##1}
      }
    \__obt_emit:
  }
\ExplSyntaxOff
\AtBeginEnvironment{longtable}{\small\setlength{\tabcolsep}{4pt}\let\mathsf\obtablesf}
\AtBeginEnvironment{tabular}{\let\mathsf\obtablesf}

% ---- Title block polish (pandoc \maketitle) ----
% Title and subtitle both centered under each other. Body metadata stays plain
% roman (never a monospaced identity dump). Public monographs must not put
% internal programme IDs (PO-xxx, GS-xxx, …) in the title line.
\usepackage{titling}
\newcommand{\obsubtitle}{}
% Capture pandoc's \subtitle{...} (provide later is a no-op once this is defined).
\newcommand{\subtitle}[1]{\gdef\obsubtitle{#1}}
\pretitle{\begin{center}\LARGE\bfseries}
\posttitle{%
  \par
  \ifx\obsubtitle\empty\else
    \vskip0.4em
    {\large\itshape \obsubtitle\par}%
  \fi
  \end{center}\vskip0.75em}
\preauthor{\begin{center}\large}\postauthor{\par\end{center}}
\predate{\begin{center}\normalsize}\postdate{\par\end{center}\vskip0.5em\hrule\vskip1em}

% ---- Status-box markers (emoji rendered as text glyphs; one place, all papers) ----
\newunicodechar{§}{\S}
\newunicodechar{–}{--}
\newunicodechar{—}{---}
\newunicodechar{‘}{`}
\newunicodechar{’}{'}
\newunicodechar{“}{``}
\newunicodechar{”}{''}
\newunicodechar{✅}{\um{\checkmark}}
\newunicodechar{⚠}{\textbf{!}}
\newunicodechar{☣}{\um{\boxtimes}}
\newunicodechar{⊣}{\um{\dashv}}
\newunicodechar{🔷}{\textcolor{oblink}{\um{\blacklozenge}}}
\newunicodechar{️}{}
\usepackage{bookmark}
\IfFileExists{xurl.sty}{\usepackage{xurl}}{} % add URL line breaks if available
\urlstyle{same}
\hypersetup{
  colorlinks=true,
  linkcolor={oblink},
  filecolor={Maroon},
  citecolor={oblink},
  urlcolor={oblink},
  pdfcreator={LaTeX via pandoc}}

\author{}
\date{2026-10-15}

\begin{document}

\section{Local Green, Global Red IV: Finite-Stage Solvability and
Completion-Relative
Obstruction}\label{local-green-global-red-iv-finite-stage-solvability-and-completion-relative-obstruction}

\textbf{Author:} Jeremy H. Carroll

\textbf{Date:} 2026-10-15

\textbf{Version:} 0.1.1

\subsection{Abstract}\label{abstract}

Evaluator agreements across finite stages of an inverse system need not
imply the existence of a solution in the original carrier. On a finite
chain, one element of the last stage pulls back along the transition
maps to a nonempty compatible thread. The general statement that a
cofiltered inverse limit of nonempty finite sets is nonempty is cited
from Stacks Tag 086J. Thread existence does not imply realization by an
element of the original carrier: in an explicit \(2\)-adic arithmetic
family, the congruence \(3x+1 \equiv 0 \pmod{2^n}\) admits an exact
solution \(x_n = -\sum_{k=0}^{n-1}(-2)^k\) at every stage \(n\) and
forms a compatible thread, yet no integer satisfies all stages
simultaneously because \(2^n \mid 3x+1\) for all \(n \ge 1\) forces
\(x = -1/3 \notin \mathbb{Z}\). The unique solution lives in the
completion \(\mathbb{Z}_2\). We formulate a general reflection image
criterion: completed solutions descend to original-carrier solutions if
and only if they belong to the image of the canonical reflection map
\(\eta\). For a finite list of families, the overlap defect \(2(B-k^*)\)
equals the dual value used by the rational certificate, and it vanishes
when a common element exists. The general statement that the
Mittag--Leffler condition implies \(R^1 \varprojlim A_n = 0\) is cited
from Stacks Tag 07KV.

\begin{center}\rule{0.5\linewidth}{0.5pt}\end{center}

\subsection{1. Problem and Relation to the LGGR
Series}\label{problem-and-relation-to-the-lggr-series}

The foundational papers in this series address evaluator composition
over a fixed finite observation presentation:

\begin{itemize}
\tightlist
\item
  \textbf{LGGR I} {[}1{]} examines sample identities, weighted
  path-kernel composition, and carrier collisions under binary evaluator
  defects.
\item
  \textbf{LGGR II} {[}2{]} establishes the organizational composition of
  boundary owners, distinguishing definitional sufficiency from true
  factorization and identifying quotient-level boundary defects.
\item
  \textbf{LGGR III} {[}3{]} layers this defect across six explicit
  system roles under finite sums and quotient descent.
\item
  \textbf{LGGR V} {[}4{]} characterizes linear observation blind spots,
  showing that full interface coverage can leave an undetected subspace
  \(\ker Q\) whenever \(\operatorname{im} D \cap \ker Q \neq \{0\}\).
\item
  \textbf{LGGR VI} {[}5{]} establishes that downstream nonlinear
  post-processing cannot recover directions already erased by an
  observation operator \(Q\).
\end{itemize}

Each of those results operates on a single, fixed stage of observation.
\textbf{LGGR IV} addresses an orthogonal mathematical question:

\begin{quote}
\emph{Can a system appear solvable across an unbounded tower of finite
observation stages, while admitting no solution in the original
carrier---even though a valid solution emerges after completion?}
\end{quote}

When an engineered system or formal evaluator is subjected to
progressively refined testing stages \(\pi_n: X \to X_n\), one observed
outcome is that every finite stage passes: for every \(n\), there exists
a state \(x_n\) satisfying the \(n\)-th requirement. If these
requirements are mutually compatible, one might hope that a global state
\(x \in X\) exists satisfying the requirements across all stages
simultaneously.

This paper establishes that such an inference is invalid in general:

\begin{enumerate}
\def\labelenumi{\arabic{enumi}.}
\tightlist
\item
  \textbf{Local Green}: Every finite refinement stage \(n\) admits a
  valid witness.
\item
  \textbf{Global Red in the Original Carrier}: No element in the
  original carrier \(X\) satisfies all stages.
\item
  \textbf{Completion Green}: A valid global witness exists in a
  completion \(\widehat{X}\).
\end{enumerate}

We show that solving a problem in the completion does not retroactively
solve the problem in the original carrier. We formalize this distinction
through a reflection image criterion and clarify the homological
boundary where lifting obstructions vanish.

\begin{center}\rule{0.5\linewidth}{0.5pt}\end{center}

\subsection{2. Finite-Stage Descent and the Moving-Witness
Criterion}\label{finite-stage-descent-and-the-moving-witness-criterion}

Let \(X\) denote a state space, and let \(I\) be a directed index set.
Consider an inverse system of quotient or observation projections:

\[
\{\pi_i: X \longrightarrow X_i\}_{i \in I}.
\]

We formalize what it means for a defect or property to factor through a
finite stage.

\subsubsection{Definition 2.1 (Finite
Descent)}\label{definition-2.1-finite-descent}

Let \(D: X \to Y\) be an evaluation or defect function. We say that
\(D\) has \textbf{finite descent} through the family
\(\{\pi_i\}_{i \in I}\), denoted \(\operatorname{FinDesc}(D)\), if there
exists an index \(i \in I\) and a function \(d_i: X_i \to Y\) such that

\[
D = d_i \circ \pi_i.
\]

When \(D = d_i \circ \pi_i\), any two states \(x, y \in X\) that agree
at stage \(i\) (that is, \(\pi_i(x) = \pi_i(y)\)) must satisfy
\(D(x) = D(y)\). The failure of finite descent is characterized by the
following criterion.

\subsubsection{Theorem 2.1 (Moving-Witness Collision
Criterion)}\label{theorem-2.1-moving-witness-collision-criterion}

Let \(X\) and \(Y\) be sets with \(Y\) nonempty, and let
\(\{\pi_i: X \to X_i\}_{i \in I}\) be a family of surjective observation
maps. A defect \(D: X \to Y\) fails to have finite descent through
\(\{\pi_i\}_{i \in I}\) if and only if for every stage \(i \in I\),
there exists a pair of states \(x_i, y_i \in X\) such that

\[
\pi_i(x_i) = \pi_i(y_i) \quad \text{and} \quad D(x_i) \neq D(y_i).
\]

\emph{Proof.}

\begin{itemize}
\tightlist
\item
  \((\implies)\) If \(\operatorname{FinDesc}(D)\) fails, then for every
  \(i \in I\), no function \(d_i: X_i \to Y\) satisfies
  \(D = d_i \circ \pi_i\). Because each \(\pi_i\) is surjective, a
  function \(d_i: X_i \to Y\) satisfying \(D = d_i \circ \pi_i\) exists
  if and only if \(D\) is constant on the fibers of \(\pi_i\) (that is,
  \(\pi_i(x) = \pi_i(y) \implies D(x) = D(y)\) for all \(x, y \in X\)).
  Without surjectivity, fiber constancy only defines \(d_i\) on the
  image \(\operatorname{im} \pi_i \subseteq X_i\), requiring an
  extension to \(X_i \setminus \operatorname{im} \pi_i\) using
  \(Y \neq \emptyset\). Under surjectivity (or extension), the failure
  of factorization implies that for every \(i \in I\), some fiber
  contains at least two points on which \(D\) takes different values:
  that is, states \(x_i, y_i \in X\) with \(\pi_i(x_i) = \pi_i(y_i)\)
  and \(D(x_i) \neq D(y_i)\).
\item
  \((\impliedby)\) Conversely, if such a pair \((x_i, y_i)\) exists for
  every \(i \in I\), then no function \(d_i: X_i \to Y\) can satisfy
  both \(d_i(\pi_i(x_i)) = D(x_i)\) and \(d_i(\pi_i(y_i)) = D(y_i)\),
  since \(\pi_i(x_i) = \pi_i(y_i)\) forces
  \(d_i(\pi_i(x_i)) = d_i(\pi_i(y_i))\) whereas \(D(x_i) \neq D(y_i)\).
  Hence no stage \(i\) factors \(D\), and \(\operatorname{FinDesc}(D)\)
  is false. \(\square\)
\end{itemize}

The Lean development formalizes this equivalence for a finite complete
state list, using one default output value to extend a stage map off its
image (\texttt{LggrIV.finite\_notFinDesc\_iff\_collisions}). It does not
treat an arbitrary non-enumerated carrier as a separate theorem.

\subsubsection{Remark 2.1 (Moving Witnesses vs.~Fixed
Collisions)}\label{remark-2.1-moving-witnesses-vs.-fixed-collisions}

The pairs \((x_i, y_i)\) depend on the stage \(i\). If the projective
family \(\{\pi_i\}_{i \in I}\) is \textbf{jointly separating}---meaning
that for every pair \(x \neq y\) in \(X\), there exists some index
\(j \in I\) such that \(\pi_j(x) \neq \pi_j(y)\)---then no \emph{fixed}
pair \((x, y)\) can satisfy \(\pi_i(x) = \pi_i(y)\) for all \(i \in I\).
Thus, failure of finite descent in a separating system is necessarily
driven by moving witnesses: as the stage \(i\) deepens, the colliding
pair \((x_i, y_i)\) must shift.

\begin{center}\rule{0.5\linewidth}{0.5pt}\end{center}

\subsection{3. A Non-Obstruction Theorem for Finite Repair
Fibers}\label{a-non-obstruction-theorem-for-finite-repair-fibers}

Before demonstrating how global solvability can fail, we establish a
structural regime where failure does not arise: namely, when the repair
sets at each finite stage are finite and nonempty.

Let \((I, \le)\) be a directed partially ordered set. An inverse system
of sets over \(I\) consists of sets \(\{S_i\}_{i \in I}\) together with
transition maps \(f_{ij}: S_j \to S_i\) for all \(i \le j\), satisfying
\(f_{ii} = \operatorname{id}_{S_i}\) and
\(f_{ik} = f_{ij} \circ f_{jk}\) for all \(i \le j \le k\).

The inverse limit is the set of compatible threads:

\[
\varprojlim_{i \in I} S_i = \left\{ (s_i)_{i \in I} \in \prod_{i \in I} S_i \;\middle|\; \forall i \le j,\; s_i = f_{ij}(s_j) \right\}.
\]

\subsubsection{Theorem 3.1 (Finite Fiber Coherence; Stacks Tag
086J)}\label{theorem-3.1-finite-fiber-coherence-stacks-tag-086j}

Let \((I, \le)\) be a directed partially ordered set, and let
\(\{S_i, f_{ij}\}\) be an inverse system of sets. If each \(S_i\) is
\textbf{finite and nonempty}, then the inverse limit is nonempty:

\[
\varprojlim_{i \in I} S_i \neq \emptyset.
\]

\emph{Proof.} Equip each finite set \(S_i\) with the discrete topology.
Because \(S_i\) is finite, it is compact and Hausdorff. By Tychonoff's
theorem, the product \(\prod_{i \in I} S_i\) is compact and Hausdorff.
For each pair \(i \le j\), the condition
\(E_{ij} = \{(s_k) \in \prod S_k \mid s_i = f_{ij}(s_j)\}\) is a closed
subset.

For any finite subset of indices \(J \subseteq I\), choose an upper
bound \(m \in I\) such that \(j \le m\) for all \(j \in J\). Since
\(S_m\) is nonempty, choose \(s_m \in S_m\) and set
\(s_j = f_{jm}(s_m)\) for \(j \in J\). This shows that any finite
intersection \(\bigcap_{i \le j \in J} E_{ij}\) is nonempty. By the
finite intersection property of compact spaces,

\[
\varprojlim_{i \in I} S_i = \bigcap_{i \le j} E_{ij} \neq \emptyset. \quad \square
\]

The Lean companion formalizes only the finite-chain case
(\texttt{LggrIV.finite\_chain\_limit\_nonempty}): one element of the
last stage, pulled back along a finite list of transition functions,
yields a nonempty compatible thread. Tag 086J remains a classical
citation, and the topological compactness argument above is cited from
the classical literature.

\subsubsection{Corollary 3.1 (Non-Obstruction for Finite Diagnostic
Repair)}\label{corollary-3.1-non-obstruction-for-finite-diagnostic-repair}

Suppose an evaluator pipeline defines for each stage \(i\) a set of
valid repairs \(R_i \subseteq X_i\). If each \(R_i\) is finite and
nonempty, and transition maps preserve validity
(\(f_{ij}(R_j) \subseteq R_i\)), then there exists a coherent global
repair sequence \((r_i)_{i \in I} \in \varprojlim R_i\).

\subsection{Consequently, an obstruction between finite-stage
solvability and global thread existence requires infinite fibers at the
finite stages, whereas an obstruction between thread existence and
original-carrier descent persists even when all finite repair sets are
singletons.}\label{consequently-an-obstruction-between-finite-stage-solvability-and-global-thread-existence-requires-infinite-fibers-at-the-finite-stages-whereas-an-obstruction-between-thread-existence-and-original-carrier-descent-persists-even-when-all-finite-repair-sets-are-singletons.}

\subsection{\texorpdfstring{4. The Explicit \(3x+1 \equiv 0 \pmod{2^n}\)
Arithmetic
Counterexample}{4. The Explicit 3x+1 \textbackslash equiv 0 \textbackslash pmod\{2\^{}n\} Arithmetic Counterexample}}\label{the-explicit-3x1-equiv-0-pmod2n-arithmetic-counterexample}

We now construct an elementary arithmetic inverse system that
illustrates the triple phenomenon:
\[\boxed{ \text{finite-stage green} \;+\; \text{original-carrier red} \;+\; \text{completion green} }\]

\subsubsection{4.1 System Definition}\label{system-definition}

Let the original carrier be the ring of integers:

\[
X = \mathbb{Z}.
\]

For each integer \(n \ge 1\), let the \(n\)-th observation stage be the
residue ring:

\[
X_n = \mathbb{Z}/2^n\mathbb{Z},
\]

with canonical projection
\(\pi_n: \mathbb{Z} \to \mathbb{Z}/2^n\mathbb{Z}\) given by
\(\pi_n(x) = x \bmod 2^n\). The transition maps
\(f_{nm}: \mathbb{Z}/2^m\mathbb{Z} \to \mathbb{Z}/2^n\mathbb{Z}\) for
\(m \ge n\) are the standard quotient projections.

The goal at each stage \(n\) is to satisfy the linear affine congruence:

\[
3x + 1 \equiv 0 \pmod{2^n}.
\]

\subsubsection{4.2 Finite-Stage Solvability (Local
Green)}\label{finite-stage-solvability-local-green}

For each \(n \ge 1\), define the integer:

\[
x_n = -\sum_{k=0}^{n-1} (-2)^k.
\]

\subsubsection{Proposition 4.1 (Exact Finite-Stage
Solution)}\label{proposition-4.1-exact-finite-stage-solution}

For every \(n \ge 1\), the integer \(x_n\) satisfies:

\[
3x_n + 1 = (-2)^n.
\]

In particular, \(2^n \mid (3x_n + 1)\), so \(\pi_n(x_n)\) is an exact
solution in \(\mathbb{Z}/2^n\mathbb{Z}\).

\emph{Proof.} Recall the geometric series identity for any real or
algebraic variable \(t\): \[
(1 - t)\sum_{k=0}^{n-1} t^k = 1 - t^n.
\] Setting \(t = -2\), we obtain \(1 - t = 1 - (-2) = 3\). Therefore: \[
3 \sum_{k=0}^{n-1} (-2)^k = 1 - (-2)^n.
\] Multiplying through by \(-1\): \[
3 \left( -\sum_{k=0}^{n-1} (-2)^k \right) = (-2)^n - 1.
\] Substituting the definition \(x_n = -\sum_{k=0}^{n-1} (-2)^k\): \[
3x_n = (-2)^n - 1 \iff 3x_n + 1 = (-2)^n.
\] Because \((-2)^n = (-1)^n 2^n\), the term \(3x_n + 1\) is an exact
multiple of \(2^n\): \[
3x_n + 1 \equiv 0 \pmod{2^n}. \quad \square
\]

\subsubsection{Proposition 4.2 (Compatibility Across
Stages)}\label{proposition-4.2-compatibility-across-stages}

The sequence of residues \((\pi_n(x_n))_{n \ge 1}\) forms a compatible
thread in \(\varprojlim \mathbb{Z}/2^n\mathbb{Z}\).

\emph{Proof.} For any \(n \ge 1\): \[
x_{n+1} - x_n = -\sum_{k=0}^n (-2)^k - \left( -\sum_{k=0}^{n-1} (-2)^k \right) = -(-2)^n.
\] Since \(-(-2)^n = (-1)^{n+1} 2^n\), we have
\(x_{n+1} \equiv x_n \pmod{2^n}\). Thus, under the projection
\(f_{n,n+1}: \mathbb{Z}/2^{n+1}\mathbb{Z} \to \mathbb{Z}/2^n\mathbb{Z}\),
\[
f_{n,n+1}(\pi_{n+1}(x_{n+1})) = \pi_n(x_n). \quad \square
\]

The first four terms of this sequence illustrate the progression:

\begin{itemize}
\tightlist
\item
  \(n = 1\): \(x_1 = -1 \implies 3(-1) + 1 = -2 \equiv 0 \pmod 2\).
\item
  \(n = 2\):
  \(x_2 = -(1 - 2) = 1 \implies 3(1) + 1 = 4 \equiv 0 \pmod 4\).
\item
  \(n = 3\):
  \(x_3 = -(1 - 2 + 4) = -3 \implies 3(-3) + 1 = -8 \equiv 0 \pmod 8\).
\item
  \(n = 4\):
  \(x_4 = -(1 - 2 + 4 - 8) = 5 \implies 3(5) + 1 = 16 \equiv 0 \pmod{16}\).
\end{itemize}

Every finite observation stage evaluates to zero defect (Local Green).

\subsubsection{4.3 Original-Carrier Failure (Global
Red)}\label{original-carrier-failure-global-red}

\subsubsection{Theorem 4.1 (Impossibility of an Integral Global
Solution)}\label{theorem-4.1-impossibility-of-an-integral-global-solution}

There exists no integer \(x \in \mathbb{Z}\) such that

\[
3x + 1 \equiv 0 \pmod{2^n} \quad \text{for all } n \ge 1.
\]

\emph{Proof.} Suppose there exists \(x \in \mathbb{Z}\) satisfying
\(3x + 1 \equiv 0 \pmod{2^n}\) for all \(n \ge 1\). This implies that
\(2^n\) divides the fixed integer \(3x + 1\) for every positive integer
\(n\):

\[
\forall n \ge 1, \quad 2^n \mid (3x + 1).
\]

In the integers, the only element divisible by \(2^n\) for all
\(n \ge 1\) is \(0\). Therefore:

\[
3x + 1 = 0 \implies x = -\frac{1}{3}.
\]

However, \(-1/3 \notin \mathbb{Z}\). Thus no such integer exists.
\(\square\)

\subsubsection{4.4 Completion Solvability (Completion
Green)}\label{completion-solvability-completion-green}

The inverse limit of the finite residue rings
\(\mathbb{Z}/2^n\mathbb{Z}\) defines the ring of \(2\)-adic integers:

\[
\mathbb{Z}_2 = \varprojlim_{n \ge 1} \mathbb{Z}/2^n\mathbb{Z}.
\]

There is a canonical injective ring homomorphism
\(\eta: \mathbb{Z} \to \mathbb{Z}_2\) mapping an integer to its coherent
sequence of residues \((\pi_n(x))_{n \ge 1}\).

\subsubsection{Classical record 4.3 (existence and uniqueness in the
2-adic
completion)}\label{classical-record-4.3-existence-and-uniqueness-in-the-2-adic-completion}

This argument is cited from standard \(2\)-adic analysis {[}9{]}.

In \(\mathbb{Z}_2\), the equation \(3x + 1 = 0\) has a unique solution:

\[
x_\infty = -\frac{1}{3} \in \mathbb{Z}_2.
\]

Moreover, the thread \((\pi_n(x_n))_{n \ge 1}\) converges \(2\)-adically
to \(x_\infty\).

\emph{Proof.} The prime \(2\) does not divide \(3\). In the ring
\(\mathbb{Z}_2\), an element is a unit if and only if its residue modulo
\(2\) is nonzero. Since \(3 \equiv 1 \pmod 2\), \(3\) is a unit in
\(\mathbb{Z}_2\). Hence the linear equation \(3x + 1 = 0\) has the
unique solution \(x = -3^{-1} = -1/3\).

The sequence of partial sums \(x_n = -\sum_{k=0}^{n-1} (-2)^k\)
satisfies: \[
|x_n - x_\infty|_2 = \left| -\frac{1 - (-2)^n}{3} - \left(-\frac{1}{3}\right) \right|_2 = \left| \frac{(-2)^n}{3} \right|_2 = \frac{1}{2^n} \xrightarrow{n \to \infty} 0.
\] Thus \(x_n \to -1/3\) in \(\mathbb{Z}_2\). \(\square\)

\begin{center}\rule{0.5\linewidth}{0.5pt}\end{center}

\subsection{5. Completion-Relative Versus Original-Carrier
Repair}\label{completion-relative-versus-original-carrier-repair}

The arithmetic example illustrates the structural hazard of completion
laundering: an evaluator observing passing checks at every stage \(n\)
concludes that the system is locally valid everywhere, and a solver
operating on the completion finds an exact limit point
\(-1/3 \in \mathbb{Z}_2\). However, if the operational system requires a
solution in \(\mathbb{Z}\), the system remains red.

\subsubsection{5.1 The Reflection Image
Criterion}\label{the-reflection-image-criterion}

Let \(X\) be an arbitrary carrier space, and let
\(\eta: X \to \widehat{X}\) denote its completion or limit embedding
with respect to an inverse family of projections.

Let \(R_X \subseteq X\) denote the predicate of valid states in the
original carrier, and let \(R_{\widehat{X}} \subseteq \widehat{X}\)
denote the extended predicate in the completion. We require that the
extension is faithful on the image of \(X\):

\[
\eta(R_X) = R_{\widehat{X}} \cap \operatorname{im}(\eta).
\]

\subsubsection{Criterion 5.1 (Reflection Descent
Criterion)}\label{criterion-5.1-reflection-descent-criterion}

Let \(\hat{r} \in R_{\widehat{X}}\) be a valid solution in the
completion. Then \(\hat{r}\) descends to a valid solution in the
original carrier \(X\) if and only if

\[
\hat{r} \in \operatorname{im}(\eta).
\]

When \(\hat{r} \notin \operatorname{im}(\eta)\), the element \(\hat{r}\)
constitutes a \textbf{completion-only witness}.

\subsubsection{5.2 The Completion
Residual}\label{the-completion-residual}

For a general completion, the primitive residual witnessing
original-carrier obstruction is given by the set membership failure:

\[
\operatorname{Res}_{\text{refl}}(\hat{r}) = \hat{r} \notin \operatorname{im}(\eta).
\]

When \(X\) and \(\widehat{X}\) possess abelian group structures and
\(\eta: X \hookrightarrow \widehat{X}\) is an injective homomorphism,
the obstruction can be packaged additively as a non-trivial coset in the
quotient group:

\[
[\hat{r}] \in \widehat{X} / \operatorname{im}(X).
\]

In our \(2\)-adic example, the element \(-1/3 \in \mathbb{Z}_2\) defines
a nonzero residual class in the quotient group
\(\mathbb{Z}_2 / \mathbb{Z}\).

\subsubsection{5.3 Operational Rule: No Completion
Laundering}\label{operational-rule-no-completion-laundering}

We state the operational discipline required to prevent invalid claim
promotion:

\begin{quote}
\textbf{Principle of No Completion Laundering}: Demonstrating that every
finite refinement stage \(n\) admits a witness, or demonstrating that a
coherent witness exists in a completion \(\widehat{X}\), does not
establish the existence of a solution in the original carrier \(X\). A
completed witness cannot be promoted to an operational certificate
unless its reflection descent \(\hat{r} \in \operatorname{im}(\eta)\) is
explicitly verified.
\end{quote}

\begin{center}\rule{0.5\linewidth}{0.5pt}\end{center}

\subsection{6. Scope and Non-Claims}\label{scope-and-non-claims}

To ensure claim hygiene, we explicitly fence the interpretive boundaries
of this paper:

\begin{enumerate}
\def\labelenumi{\arabic{enumi}.}
\tightlist
\item
  \textbf{Diagnostic Geometry, Not Physics or Ontology}: The divergence
  between \(\mathbb{Z}\) and \(\mathbb{Z}_2\) is a property of
  projective limits and metric completions. It does not model physical
  spacetime quantization, ontological emergence, or quantum measurement.
\item
  \textbf{Representation Limits}: The existence of moving-witness
  collisions under \(\neg \operatorname{FinDesc}(D)\) reflects the
  lossiness of finite stages, not an inherent defect in the underlying
  mathematical reality.
\item
  \textbf{No Continuum Extrapolation}: Passing a countable sequence of
  finite tests provides evidence only for those tests. It does not
  authorize unverified assertions about continuous limits.
\end{enumerate}

\begin{center}\rule{0.5\linewidth}{0.5pt}\end{center}

\subsection{Appendix: Derived Limits and the Mittag--Leffler Vanishing
Theorem}\label{appendix-derived-limits-and-the-mittagleffler-vanishing-theorem}

In this appendix, we place the lifting obstruction into the standard
framework of homological algebra. We clarify the role of the
Mittag--Leffler condition, correcting a common misconception regarding
the derived inverse limit \(R^1 \varprojlim\).

\subsubsection{A.1 The Inverse System Exact
Sequence}\label{a.1-the-inverse-system-exact-sequence}

Let \(0 \to A_n \xrightarrow{\iota_n} B_n \xrightarrow{p_n} C_n \to 0\)
be a sequence of short exact sequences of abelian groups, compatible
with transition maps:

\begin{Shaded}
\begin{Highlighting}[]
\NormalTok{0 {-}{-}{-}\textgreater{} A\_\{n+1\} {-}{-}{-}\textgreater{} B\_\{n+1\} {-}{-}{-}\textgreater{} C\_\{n+1\} {-}{-}{-}\textgreater{} 0}
\NormalTok{         |              |              |}
\NormalTok{         v              v              v}
\NormalTok{0 {-}{-}{-}\textgreater{}   A\_n   {-}{-}{-}\textgreater{}   B\_n   {-}{-}{-}\textgreater{}   C\_n   {-}{-}{-}\textgreater{} 0}
\end{Highlighting}
\end{Shaded}

Applying the left-exact inverse limit functor \(\varprojlim\) yields the
standard six-term exact sequence:

\[
0 \longrightarrow \varprojlim A_n \longrightarrow \varprojlim B_n \longrightarrow \varprojlim C_n \xrightarrow{\;\delta\;} R^1 \varprojlim A_n \longrightarrow R^1 \varprojlim B_n \longrightarrow R^1 \varprojlim C_n \longrightarrow 0.
\]

Here, \(R^1 \varprojlim A_n\) denotes the first derived limit of the
inverse system \(\{A_n\}\).

\subsubsection{A.2 The Lifting Obstruction
Class}\label{a.2-the-lifting-obstruction-class}

Suppose we are given a compatible sequence of target requirements
\((c_n) \in \varprojlim C_n\).

\begin{itemize}
\tightlist
\item
  At each finite stage \(n\), surjectivity of \(p_n\) ensures that there
  exists some \(b_n \in B_n\) lifting \(c_n\) (Finite-Stage Green).
\item
  However, the choices of \(b_n\) may not be compatible under transition
  maps: \(f_{n,n+1}(b_{n+1})\) need not equal \(b_n\).
\item
  The connecting homomorphism
  \(\delta: \varprojlim C_n \to R^1 \varprojlim A_n\) maps the
  compatible family \((c_n)\) to an obstruction class:
\end{itemize}

\[
\omega = \delta((c_n)) \in R^1 \varprojlim A_n.
\]

By exactness of the sequence, the compatible family \((c_n)\) lifts to a
global coherent element \((b_n) \in \varprojlim B_n\) if and only if
\(\omega = 0\).

\subsubsection{A.3 The Mittag--Leffler Condition (Stacks Tag 07KV, Tag
0597)}\label{a.3-the-mittagleffler-condition-stacks-tag-07kv-tag-0597}

\subsubsection{Definition A.1 (Mittag--Leffler
Condition)}\label{definition-a.1-mittagleffler-condition}

An inverse system of abelian groups \(\{A_n, f_{nm}\}\) satisfies the
\textbf{Mittag--Leffler (ML) condition} if for every \(n\), the
descending chain of images

\[
\operatorname{im}(f_{n, n+k}) \subseteq A_n \quad (k \ge 0)
\]

eventually stabilizes: that is, there exists \(k_0\) such that for all
\(k \ge k_0\),
\(\operatorname{im}(f_{n, n+k}) = \operatorname{im}(f_{n, n+k_0})\).

\subsubsection{Theorem A.1 (Mittag--Leffler Vanishing; Stacks Tag
07KV)}\label{theorem-a.1-mittagleffler-vanishing-stacks-tag-07kv}

If \(\{A_n\}\) is an inverse system of abelian groups indexed by the
natural numbers that satisfies the Mittag--Leffler condition, then:

\[
R^1 \varprojlim A_n = 0.
\]

\emph{Proof.} See Stacks Project, Tag 07KV. \(\square\)

The Lean companion formalizes only the finite overlap identity
(\texttt{LggrIV.overlap\_dual\_residual}): for a finite list of families
the primal defect \(2(B-k^*)\) equals the dual value
\(\mu - \sum \alpha_i\), and a nonempty list with \(k^*=B\) has a common
element. Tag 07KV remains a classical citation.

\subsubsection{Corollary A.1 (Vanishing of the Lifting
Obstruction)}\label{corollary-a.1-vanishing-of-the-lifting-obstruction}

If the kernel system \(\{A_n\}\) satisfies the Mittag--Leffler
condition, then every compatible family \((c_n) \in \varprojlim C_n\)
lifts to a compatible family \((b_n) \in \varprojlim B_n\).

\subsubsection{Remark A.1 (Correction to the Role of
Mittag--Leffler)}\label{remark-a.1-correction-to-the-role-of-mittagleffler}

The Mittag--Leffler condition does \emph{not} create an obstruction;
rather, it ensures the \emph{absence} of a lifting obstruction
(\(R^1 \varprojlim A_n = 0\)). Consequently, coherent lifting failure
(Local Green, Global Red in an abelian tower) can occur only when the
kernel transition system \(\{A_n\}\) fails the Mittag--Leffler
condition.

\subsubsection{\texorpdfstring{A.4 The Direct Arithmetic Bridge:
\(R^1 \varprojlim (2^n \mathbb{Z}) \cong \mathbb{Z}_2 / \mathbb{Z}\)}{A.4 The Direct Arithmetic Bridge: R\^{}1 \textbackslash varprojlim (2\^{}n \textbackslash mathbb\{Z\}) \textbackslash cong \textbackslash mathbb\{Z\}\_2 / \textbackslash mathbb\{Z\}}}\label{a.4-the-direct-arithmetic-bridge-r1-varprojlim-2n-mathbbz-cong-mathbbz_2-mathbbz}

The derived-limit exact sequence of Section A.1 applies directly to the
paper's primary arithmetic system. For each \(n \ge 1\), consider the
short exact sequence of abelian groups:

\[
0 \longrightarrow 2^n \mathbb{Z} \longrightarrow \mathbb{Z} \longrightarrow \mathbb{Z}/2^n \mathbb{Z} \longrightarrow 0,
\]

equipped with standard inclusion homomorphisms on the submodules
\(2^n \mathbb{Z}\), identity transitions on the constant tower
\(\mathbb{Z}\), and modular reductions on the residue rings
\(\mathbb{Z}/2^n \mathbb{Z}\).

Evaluating the projective limits: 1.
\(\varprojlim 2^n \mathbb{Z} = \bigcap_{n \ge 1} 2^n \mathbb{Z} = \{0\}\).
2. The constant tower \(\{\mathbb{Z}, \operatorname{id}\}\) satisfies
the Mittag--Leffler condition (its transition images are identically
\(\mathbb{Z}\)), hence \(\varprojlim \mathbb{Z} = \mathbb{Z}\) and
\(R^1 \varprojlim \mathbb{Z} = 0\). 3.
\(\varprojlim \mathbb{Z}/2^n \mathbb{Z} = \mathbb{Z}_2\) (the ring of
\(2\)-adic integers).

Substituting these evaluations into the six-term derived limit exact
sequence yields:

\[
0 \longrightarrow 0 \longrightarrow \mathbb{Z} \xrightarrow{\;\eta\;} \mathbb{Z}_2 \xrightarrow{\;\delta\;} R^1 \varprojlim (2^n \mathbb{Z}) \longrightarrow 0.
\]

Consequently, as abelian groups:

\[
R^1 \varprojlim (2^n \mathbb{Z}) \cong \mathbb{Z}_2 / \mathbb{Z}.
\]

The connecting homomorphism
\(\delta: \mathbb{Z}_2 \to R^1 \varprojlim (2^n \mathbb{Z})\) is the
canonical quotient projection modulo the embedded integers
\(\eta(\mathbb{Z})\).

In our arithmetic system, the unique \(2\)-adic solution
\(x_\infty = -1/3 \in \mathbb{Z}_2\) does not lie in the image of
\(\mathbb{Z}\). Therefore:

\[
\delta(x_\infty) = [-1/3] \neq 0 \in R^1 \varprojlim (2^n \mathbb{Z}).
\]

Thus, the completion residual \(\hat{r} \notin \operatorname{im}(\eta)\)
and the non-vanishing derived limit lifting obstruction
\(\delta(x_\infty) \neq 0\) coincide on the paper's primary system.

At the level of finite representatives, the chosen stage solutions
\((x_n)_{n \ge 1}\) satisfy the step discrepancies: \[
x_n - x_{n+1} = (-2)^n \in 2^n \mathbb{Z}.
\] If these step discrepancies were trivialized by compatible kernel
corrections \(k_n \in 2^n \mathbb{Z}\), the corrected sequence
\(x_n' = x_n - k_n\) would stabilize to a single integer
\(x \in \mathbb{Z}\) satisfying \(3x + 1 \equiv 0 \pmod{2^n}\) for all
\(n \ge 1\), which is ruled out by Theorem 4.1.

\emph{Note on quotient topology:} The isomorphism
\(R^1 \varprojlim (2^n \mathbb{Z}) \cong \mathbb{Z}_2 / \mathbb{Z}\) is
an algebraic group isomorphism. Because \(\mathbb{Z}\) is dense in
\(\mathbb{Z}_2\) under the \(2\)-adic metric, the quotient does not
carry a Hausdorff metric; non-zero cosets admit representatives
arbitrarily close to \(0\) in metric distance. This separation
reinforces Criterion 5.1: finite-stage or metric proximity cannot
substitute for exact reflection descent into the original operational
carrier.

\begin{center}\rule{0.5\linewidth}{0.5pt}\end{center}

\subsection{References}\label{references}

\begin{enumerate}
\def\labelenumi{\arabic{enumi}.}
\tightlist
\item
  Carroll, J. H. \emph{Local Green, Global Red I: Evaluated Path Kernels
  for Composed Agent Reliability.} Zenodo preprint, 2026. DOI:
  \url{https://doi.org/10.5281/zenodo.22881911}.
\item
  Carroll, J. H. \emph{Local Green, Global Red II: Organizational
  Evaluator Gluing.} Zenodo preprint, 2026. DOI:
  \url{https://doi.org/10.5281/zenodo.22882042}.
\item
  Carroll, J. H. \emph{Local Green, Global Red III: A finite evaluator
  overlay across six system roles.} Zenodo preprint, 2026.
\item
  Carroll, J. H. \emph{Local Green, Global Red V: Linear Observation
  Blind Spots and Kernel Intersections.} Zenodo preprint, 2026.
\item
  Carroll, J. H. \emph{Local Green, Global Red VI: Nonlinear Aggregation
  and Adjoint Probing.} Zenodo preprint, 2026.
\item
  The Stacks Project Authors. \emph{Tag 086J: Inverse limits of finite
  nonempty sets.} \url{https://stacks.math.columbia.edu/tag/086J}.
\item
  The Stacks Project Authors. \emph{Tag 07KV: Derived inverse limits and
  the Mittag--Leffler condition.}
  \url{https://stacks.math.columbia.edu/tag/07KV}.
\item
  The Stacks Project Authors. \emph{Tag 0597: Mittag--Leffler systems of
  sets.} \url{https://stacks.math.columbia.edu/tag/0597}.
\item
  Koblitz, N. \emph{p-adic Numbers, p-adic Analysis, and
  Zeta-Functions.} Graduate Texts in Mathematics, vol.~58, Springer,
  1984.
\item
  Weibel, C. A. \emph{An Introduction to Homological Algebra.} Cambridge
  Studies in Advanced Mathematics, vol.~38, Cambridge University Press,
  1994.
\end{enumerate}

\end{document}
